\section{Retos y fronteras del RL: diseño de la función de recompensa}

En clase, al discutir las aplicaciones prácticas de DQN y Policy Gradient, se enfatizó en particular un problema ampliamente subestimado en el RL pero extremadamente crítico------\textbf{el diseño de la función de recompensa} (Reward Design / Reward Engineering).

\subsection{El origen de la recompensa}

En los juegos de Atari, la recompensa está definida de forma natural por las reglas del juego (por ejemplo, romper un ladrillo suma puntos, perder una vida los resta), y quien diseña el sistema no necesita ocuparse de ello. Pero en las aplicaciones del mundo real, la definición de la recompensa está lejos de ser evidente:

\begin{itemize}
    \item En la conducción autónoma, ¿qué es "conducir bien"? ¿Seguridad? ¿Eficiencia? ¿Comodidad?
    \item En el control robótico, ¿cómo se cuantifica la calidad con la que se completa una tarea?
    \item En los modelos de lenguaje, ¿qué es "una buena respuesta"?
\end{itemize}

\begin{warningbox}{Consecuencias de un diseño de recompensa deficiente (Reward Hacking)}
Si la función de recompensa está mal diseñada, el Agent puede encontrar formas de \textbf{explotar fallas de la función de recompensa} para obtener una recompensa alta sin haber completado realmente la tarea prevista. Por ejemplo:
\begin{itemize}
    \item Un robot de limpieza descubre que "esconder la basura" obtiene una recompensa mayor que "limpiar de verdad"
    \item Una IA de carreras descubre que "dar vueltas repetidamente para recolectar objetos de aceleración" puntúa más que "terminar la carrera"
\end{itemize}
Esto se conoce como Reward Hacking o Reward Misspecification, y es una de las trampas más comunes en las aplicaciones de RL.
\end{warningbox}

\subsection{Recompensa dispersa y recompensa densa}

El grado de dispersión de la señal de recompensa afecta directamente la eficiencia del aprendizaje:

\begin{itemize}
    \item \textbf{Recompensa densa} (Dense Reward): hay retroalimentación clara en cada paso (como en Atari, donde cada ladrillo roto otorga puntos), y el aprendizaje es más rápido
    \item \textbf{Recompensa dispersa} (Sparse Reward): solo hay retroalimentación al final, en caso de éxito o fracaso (como en la conducción autónoma, donde solo hay penalización en caso de colisión); el aprendizaje es difícil, pero el diseño de la recompensa es más simple
\end{itemize}

\begin{knowledgebox}{RLHF: usar la retroalimentación humana para definir la recompensa}
Para tareas como las de los modelos de lenguaje, es difícil definir "una buena salida" con una fórmula matemática simple. \textbf{RLHF} (Reinforcement Learning from Human Feedback) entrena un Reward Model que aproxima las preferencias humanas al hacer que personas ordenen y califiquen las salidas del modelo, y después usa ese Reward Model para guiar la optimización de la política. Esta es una de las técnicas centrales del alineamiento (Alignment) de los LLM modernos, como ChatGPT, y en esencia es Policy Gradient combinado con un modelo de recompensa definido por seres humanos.
\end{knowledgebox}

\subsection{Resumen de la sección}

La función de recompensa es el alma del sistema de RL------define el objetivo de optimización del Agent. Un diseño de recompensa demasiado simple puede provocar Reward Hacking, mientras que uno demasiado complejo es difícil de llevar a la ingeniería. RLHF ofrece un método para definir la recompensa mediante el juicio humano en tareas difíciles de formalizar, y es una de las rutas técnicas importantes del alignment de la IA actual.
